\section{Effective bounds on the total extreme count}\label{sec:counts}

Finite local stationarity constraints give upper bounds on the
arithmetical complexity of counting extreme translation-invariant stationary states.
The argument applies to rational classical PCA and algebraically
presented local quantum circuits.

\begin{theorem}\label{thm:counts}
Uniformly on valid finite presentations, for every fixed integer
\(r\ge1\), the condition \(N\le r\) is \(\Pi^0_2\).
Consequently \(N\ge k\) is \(\Sigma^0_2\) for \(k\ge2\), and
\(N=1\) is \(\Pi^0_2\).  For \(k\ge2\), the condition \(N=k\)
belongs to \(D_2(\Sigma^0_2)\subseteq\Delta^0_3\).
Finite total extreme count is \(\Sigma^0_3\), and infinite total
extreme count is \(\Pi^0_3\).
\end{theorem}

Here \(D_2(\Sigma^0_2)\) denotes differences of two
\(\Sigma^0_2\) sets.  All the assertions are upper bounds on the
total count \(N\).

\begin{proof}
\emph{Linear descriptions of the stationary set.}
Choose an effective countable determining family of real local
coordinates \(x_1,x_2,\ldots\), each bounded in absolute value by
one.  Classically use cylinder probabilities.  Quantumly use real
and imaginary entries of finite-region density matrices, with
Hermiticity built into the coordinates.  A state is specified by
countably many linear equalities and inequalities: normalization,
consistency under restriction, translation invariance, stationarity
and positivity.  Classical positivity is entrywise nonnegativity.
For a finite quantum region \(\Lambda\), impose
\[
 v^*\rho_{\Lambda}v\ge0\qquad
 \text{for every vector }v\text{ with rational real and imaginary parts}.
\]
Each is a linear inequality in the chosen real coordinates.  Density
of these vectors makes the countable family equivalent to positive
semidefiniteness of \(\rho_{\Lambda}\).  Stationarity on a local Hermitian
basis has exact algebraic coefficients, computed from the finite
backward circuit.  Thus all constraints can be enumerated effectively
and their coefficients admit exact arithmetic and order tests.

Let \(b(m)\ge m\) increase so that the first \(m\) constraints use
only the first \(b(m)\) coordinates.  In those coordinates let
\(P_m\) be the bounded algebraic polytope defined by the first
\(m\) constraints and the coordinate bounds \([-1,1]\).
Choose the enumeration so every earlier constraint is retained.
The restriction of every genuine stationary state belongs to every
\(P_m\), so these polytopes are nonempty.  Conversely, if
\(m_j\to\infty\) and points of \(P_{m_j}\), extended arbitrarily
in the unused bounded coordinates, converge coordinatewise, their
limit satisfies every constraint.  It is therefore a genuine
translation-invariant stationary law or state.

Write \(\mathcal X=\cK(M)\).  For every fixed \(n\), the first-\(n\)
coordinate projections of \(P_m\) converge in Hausdorff distance to
the corresponding projection of \(\mathcal X\).  Indeed, a sequence of
outer points staying a positive distance away would have a
coordinatewise convergent subsequence, whose limit would belong to
\(\mathcal X\).  Genuine points are contained at every stage.  

\emph{Finite decidable approximation tests.}
The vertices of each \(P_m\) are computable with exact algebraic
arithmetic.  Enumerate full-rank systems of active constraints,
solve their linear equations and check all inequalities.  This also
covers lower-dimensional polytopes: if the active normals at a
vertex failed to span the ambient coordinate space, a small segment
in a common null direction would remain feasible.  Every bounded
polytope is the convex hull of its vertices, as follows by expressing
an interior point on a boundary-to-boundary segment and inducting on
face dimension.  Finite algebraic linear feasibility is consequently
decidable by the same finite linear-algebra operations.

For \(m\ge n\), put \(\eta=2^{-n}\), and let \(V\ge1\) be the
number of vertices of \(P_m\).  Choose an integer
\(L_{m,n}\ge8V/\eta\).  Let \(\cQ_{m,n}\) contain all convex
combinations of these vertices with nonnegative weights in
multiples of \(1/L_{m,n}\).  It is a finite subset of \(P_m\).
It is an \(\eta/4\)-net in the first \(n\) coordinates: round
the first \(V-1\) weights down to multiples of \(1/L_{m,n}\), putting
the remainder on the last vertex.  The sum of absolute weight
errors is at most \(2(V-1)/L_{m,n}\), so each coordinate error is at most
\(\eta/4\).  The case \(V=1\) is immediate.

Define the finite test \(\mathsf{Test}_r(n,m)\) to hold if some \(r\)-tuple
of points in \(\cQ_{m,n}\), with repetitions allowed, has the
following property: every vertex of \(P_m\) is within \(\eta\),
in first-\(n\) coordinate maximum norm, of that tuple's convex
hull.  For a fixed tuple and vertex \(v\), the test is feasibility
of
\[
 \lambda_i\ge0,\qquad \sum_{i=1}^r\lambda_i=1,\qquad
 \left|v_j-\sum_{i=1}^r\lambda_i y_{i,j}\right|\le\eta
 \quad(1\le j\le n).
\]
This is a finite algebraic linear-feasibility problem.  Define the test to be false on malformed or empty stage data.
Enumerating the finitely many tuples makes \(\mathsf{Test}_r(n,m)\) decidable.  Covering
the vertices also covers \(P_m\), since the sum of a convex hull
and a maximum-norm ball is convex.

\emph{The quantifier characterization.}
We claim that
\begin{equation}\label{eq:finite-count-formula}
 N\le r
 \quad\Longleftrightarrow\quad
 \forall n\ge1\ \exists m\ge n:\ \mathsf{Test}_r(n,m).
\end{equation}
First recall why a compact convex set with finitely many extreme
points is their convex hull in this setting.  Suppose a point of
\(\mathcal X\) lies outside the hull of the alleged finite list.  Compactness
of the coefficient simplex gives a finite coordinate projection
in which it still lies outside that hull.  Finite-dimensional
separation then supplies a local linear functional whose maximum
on \(\mathcal X\) exceeds its maximum on the list.  Its maximizing compact
face is disjoint from the hull.  Successively maximizing a
countable determining coordinate family inside that face leaves a
singleton intersection.  This point is extreme in \(\mathcal X\),
contradicting completeness of the list.

If \(N\le r\), choose \(r\) generators of \(\mathcal X\), repeating
points when necessary.  At fixed \(n\), choose \(m\ge n\) so
large that the first-\(n\) projection of \(P_m\) lies within
\(\eta/4\) of that of \(\mathcal X\).  Approximate the \(r\) generators
by points of \(\cQ_{m,n}\) within \(\eta/4\) in those
coordinates.  Their hull then covers the projected \(P_m\) within
\(\eta/2\), so \(\mathsf{Test}_r(n,m)\) holds.

Conversely, choose stages \(m_n\ge n\) and successful centers
\(y_1^{(n)},\ldots,y_r^{(n)}\in P_{m_n}\).
After extending unused coordinates in \([-1,1]\), simultaneous
product compactness gives a subsequence along which every center
converges coordinatewise to a fixed point \(y_i\in \mathcal X\).
Membership in \(\mathcal X\) follows because \(m_n\ge n\to\infty\).
Now fix any \(x\in \mathcal X\).  Its restriction lies in each relevant
\(P_{m_n}\), so vertex coverage supplies simplex coefficients
\(\lambda_i^{(n)}\) approximating its first \(n\) coordinates
within \(2^{-n}\).  A further subsequence makes these coefficients
converge.  Passing to each fixed coordinate yields
\[
 x=\sum_{i=1}^r\lambda_i y_i.
\]
The limiting centers were chosen before \(x\); only the
coefficients depend on it.  Hence \(\mathcal X\) is the convex hull of
these same \(r\) points and has at most \(r\) extreme points.
This proves~\eqref{eq:finite-count-formula}.

The right side is a universal-existential formula with decidable
matrix, proving the \(\Pi^0_2\) bound.  Its complement gives
\(N\ge k\in\Sigma^0_2\) for \(k\ge2\).  Since \(\mathcal X\) is
nonempty compact convex, it has an extreme point, so \(N=1\) is
already the case \(N\le1\).  For \(k\ge2\), write
\[
 \{N=k\}=\{N\ge k\}\setminus\{N\ge k+1\}.
\]
This gives \(D_2(\Sigma^0_2)\), contained in \(\Delta^0_3\).
Finally, prefixing~\eqref{eq:finite-count-formula} with an
existential integer \(r\) makes finite total count \(\Sigma^0_3\);
its complement is \(\Pi^0_3\).
\end{proof}




